
At the energy scales relevant for the dynamics of atomic nuclei, the fundamental theory of strong interactions, quantum chromodynamics (QCD), becomes non-perturbative, and quarks and gluons are confined within hadrons. Following the seminal works of Weinberg~\cite{weinberg_nuclear_1990,weinberg_effective_1991}, effective field theories (EFTs) have emerged as the framework of choice for computing low-energy nuclear observables with quantified uncertainties. Two prominent examples are pionless EFT and chiral EFT, both of which exploit a separation between a ``hard'' momentum scale $\Lambda$ and a ``soft'' scale $Q$ that characterizes the typical momenta in low-energy nuclear systems~\cite{epelbaum_modern_2009,machleidt_chiral_2011}. Pionless EFT is applicable when $Q \ll m_\pi$, and its breakdown scale is set by $\Lambda \sim m_\pi$. In contrast, chiral EFT retains explicit pion degrees of freedom, and its breakdown scale is considerably higher, $\Lambda \sim 600$--$700$ MeV, while the soft scale remains $Q \sim m_\pi$. At low energies, protons and neutrons interact through nonrelativistic potentials systematically organized as an expansion in $Q/\Lambda$~\cite{van_kolck_few_1994,epelbaum_modern_2009,machleidt_chiral_2011}. High-momentum physics is integrated out and encoded in low-energy constants fitted to nucleon--nucleon scattering data and observables in light nuclei. 

Interactions and electroweak currents derived consistently within the EFT framework provide the primary input to \emph{ab initio} methods designed to solve the nuclear many-body Schr\"odinger equation with controlled and systematically improvable approximations~\cite{hergert_guided_2020,ekstrom_what_2023}. Single-particle basis methods such as the no-core shell model (NCSM)~\cite{barrett_ab_2013}, coupled-cluster (CC) theory~\cite{hagen_coupled-cluster_2014}, the in-medium similarity renormalization group (IMSRG)~\cite{hergert_-medium_2016}, and self-consistent Green's function (SCGF) techniques~\cite{Carbone:2013eqa} have achieved impressive success. The NCSM and its resonating-group extension~\cite{navratil_recent_2009} accurately describe clustering and collective dynamics~\cite{quaglioni_ab_2008,roth_similarity-transformed_2011,mccoy_emergent_2020} but are generally restricted to light nuclei, and convergence for long-range observables can be slow~\cite{caprio_robust_2022}. Polynomially scaling approaches such as CC and IMSRG now reach nuclei as heavy as $^{100}$Sn~\cite{hagen_structure_2016,morris_structure_2018,mougeot_mass_2021} and $^{208}$Pb~\cite{hu_ab_2022}, providing access to spectra, response functions, and transition rates~\cite{sobczyk_ab_2021,acharya_16o_2024,bonaiti_electromagnetic_2024}. These developments have deepened our understanding of the long-standing quenching of $g_A$~\cite{gysbers_discrepancy_2019} and enabled uncertainty-quantified predictions of neutrinoless double-$\beta$ decay~\cite{yao_ab_2020,wirth_ab_2021,belley_ab_2021}. Nevertheless, truncation of model spaces limits their ability to fully capture multi-scale phenomena, especially alpha clustering and short-range correlations. Symmetry-adapted shell-model approaches~\cite{launey_symmetry-guided_2016,dytrych_physics_2020} mitigate some of these challenges but they may not  be able to fully capture short-distance dynamics.

A second class of \emph{ab initio} methods relies on stochastic formulations, such as nuclear lattice EFT (NLEFT)~\cite{lee_lattice_2025} and continuum quantum Monte Carlo (QMC) techniques~\cite{carlson_quantum_2015,lynn2019,gandolfi2020}. NLEFT has advanced to medium-mass nuclei~\cite{elhatisari_wavefunction_2024}, revealing compelling signatures of alpha clustering~\cite{shen_emergent_2023} and providing insights into nuclear matter at finite temperature~\cite{Ren:2023ued,ma_structure_2024}, enabled by substantial algorithmic progress~\cite{sarkar_floating_2023}. However, the fermion sign problem and the computational cost of fine lattice spacings limit its ability to resolve short-range dynamics with high accuracy. Continuum QMC methods, such as Green’s function Monte Carlo (GFMC)~\cite{carlson_alpha_1988} and auxiliary-field diffusion Monte Carlo (AFDMC)~\cite{Schmidt:1999lik}, accurately describe nuclear dynamics across long, intermediate, and short ranges. In practice, however, GFMC scales exponentially with the mass number, restricting applications to nuclei with $A \lesssim 13$, while AFDMC encounters a severe sign problem when realistic correlations are introduced, limiting its applicability to nuclei with $A \sim 16$~\cite{novario_trends_2023,martin_auxiliary_2023}.



These challenges have motivated the exploration of neural-network representations of quantum many-body wave functions, a framework that has already achieved notable success in condensed-matter physics and quantum chemistry~\cite{carleo_solving_2017,hermann_deep-neural-network_2020,pfau_ab_2020}. In the nuclear domain, however, neural-network quantum states (NQS) methods face a distinct set of difficulties. Nuclear interactions are highly non-perturbative and exhibit strong spin--isospin dependence, placing them well outside the regime where mean-field approximations provide a reliable starting point~\cite{wiringa_accurate_1995}. As a result, NQS methods are confronted with correlations at all length scales. At the same time, realistic nuclear Hamiltonians require treating both continuous spatial degrees of freedom and discrete spin and isospin degrees of freedom, significantly enlarging the structure of the underlying Hilbert space.

In this review, we discuss the development and application of NQS ans\"atze in nuclear physics and closely related areas. Following the pioneering application of NQS to the deuteron~\cite{keeble_machine_2020}, variational Monte Carlo methods based on NQS~\cite{adams_variational_2021,yang_consistent_2022,lovato_hidden-nucleons_2022,Yang:2023rlw,gnech_distilling_2024} have proven to provide a systematically improvable framework with polynomial computational scaling and have begun to address several long-standing challenges in \emph{ab initio} nuclear theory. Our focus is on first-quantized architectures, which operate directly on the continuous spatial coordinates and discrete spin--isospin degrees of freedom of the nucleons. This representation naturally accommodates short-, intermediate-, and long-range nuclear dynamics, from the repulsive core and tensor correlations to clustering and collective degrees of freedom. Moreover, NQS approaches capture the multiscale behavior characteristic of nuclei and neutron-star matter~\cite{fore2023}, including superfluidity, the emergence of nuclear clusters, and dominant short-range correlations~\cite{kim_neural-network_2024,fore_investigating_2024}.

First, we briefly review existing continuum QMC methods, focusing in particular on GFMC and the computational challenges associated with treating nuclei with more than $A=13$ nucleons. This naturally leads to a discussion on how the Hilbert space is sampled in VMC methods, with emphasis on the nuclear case, where both continuous (spatial coordinates) and discrete (spin--isospin) degrees of freedom must be sampled efficiently. After providing a short recap of the elements of deep learning relevant to NQS approaches, we introduce the main first-quantized NQS architectures proposed to date and discuss how physical symmetries—such as parity, time-reversal, and translation invariance—are incorporated into these models.

In the second part of the review, we summarize the most important recent results for finite nuclei and for infinite nuclear and neutron matter. We then discuss initial applications of NQS wave functions to lepton--nucleus and nucleus--nucleus scattering, highlighting the potential of these methods to deliver unified and accurate \emph{ab initio} descriptions of both nuclear structure and reactions, and to extend the reach of nuclear many-body theory beyond the traditional limits imposed by basis truncations and the sign problem. Before concluding, we provide a selective overview NQS approaches to condensed matter theory, focusing particularly on those that have direct links to nuclear theory.

